\documentclass[11pt,a4paper]{article}

% --- 1. Configuración de Página y Márgenes ---
\usepackage[margin=1in]{geometry}
\usepackage{authblk}

% --- 2. Paquetes Matemáticos ---
\usepackage{amsmath, amssymb, amsthm}
\usepackage{mathtools}
\usepackage{bm}
\usepackage{dsfont}
\usepackage{physics}
\usepackage{braket}

% --- 3. Elementos Gráficos y Tablas ---
\usepackage{graphicx}
\usepackage{color}
\usepackage{xcolor}
\usepackage{booktabs}

% --- 4. Algoritmos y Código ---
\usepackage{algorithm}
\usepackage{algpseudocode}
\usepackage{listings}

% --- Listas ---
\usepackage[shortlabels]{enumitem}

% --- 5. Hipervínculos ---
\usepackage[numbers,sort&compress]{natbib}
\usepackage[colorlinks=true, linkcolor=blue!70!black, citecolor=red!70!black, urlcolor=blue!70!black]{hyperref}
\usepackage{orcidlink}
\newcommand{\orcid}[1]{\orcidlink{#1}}
\usepackage{url}

% --- 6. Cleveref ---
\usepackage{cleveref}
\crefname{equation}{Eq.}{Eqs.}
\crefname{section}{Sec.}{Secs.}
\crefname{subsection}{Sec.}{Secs.}
\crefname{table}{Table}{Tables}
\crefname{figure}{Fig.}{Figs.}
\labelformat{lemma}{Lemma~#1}
\labelformat{theorem}{Theorem~#1}

% --- Estilos de Código ---
\definecolor{codegreen}{rgb}{0,0.6,0}
\definecolor{codegray}{rgb}{0.5,0.5,0.5}
\definecolor{codepurple}{rgb}{0.58,0,0.82}
\definecolor{backcolour}{rgb}{0.95,0.95,0.92}

\lstdefinestyle{mystyle}{
    backgroundcolor=\color{backcolour},
    commentstyle=\color{codegreen},
    keywordstyle=\color{magenta},
    numberstyle=\tiny\color{codegray},
    stringstyle=\color{codepurple},
    basicstyle=\ttfamily\footnotesize,
    breakatwhitespace=false,
    breaklines=true,
    captionpos=b,
    keepspaces=true,
    numbers=left,
    numbersep=5pt,
    showspaces=false,
    showstringspaces=false,
    showtabs=false,
    tabsize=2,
    language=Python
}
\lstset{style=mystyle}

% --- Entornos de teoremas ---
\newtheorem{theorem}{Theorem}[section]
\newtheorem{lemma}[theorem]{Lemma}
\newtheorem{corollary}[theorem]{Corollary}
\newtheorem{definition}[theorem]{Definition}
\newtheorem{remark}[theorem]{Remark}
\newtheorem{proposition}[theorem]{Proposition}

% --- Macros ---
\newcommand{\Z}{\mathbb{Z}}
\newcommand{\C}{\mathbb{C}}
\newcommand{\R}{\mathbb{R}}
\newcommand{\N}{\mathbb{N}}
\newcommand{\Hilb}{\mathcal{H}}
\newcommand{\Zmod}{\mathbb{Z}/6\mathbb{Z}}
\DeclareMathOperator{\IPR}{IPR}
\DeclareMathOperator{\Var}{Var}
\DeclareMathOperator{\E}{\mathbb{E}}

\begin{document}

\title{\textbf{First-neighbor connectivity governs the non-ergodic extended phase in power-law random banded matrices with sparse hopping constraints}}

\author{ José Ignacio Peinador Sala \orcidlink{0009-0008-1822-3452}}
\affil{Independent Researcher, Valladolid, Spain\\ \texttt{joseignacio.peinador@gmail.com}}

\date{\today}

\maketitle

\begin{abstract}
We investigate the non-ergodic extended (NEE) phase in power-law random
banded matrices (PRBM) subjected to sparse arithmetic hopping constraints.
The diagonal potential deterministically reproduces the Riemann zeta zeros'
macroscopic density via the Lambert $W$ function, and the off-diagonal
disorder is constrained by masks eliminating specific hopping channels.
Exact diagonalization up to $N=20\,000$ in double precision reveals three
mutually consistent signatures of the NEE phase for the coprime-to-6 mask:
(i) a sub-diffusive spectral form factor ramp with $\gamma \approx 0.61$;
(ii) a Chalker dynamical scaling relation $\beta = 2D_2$ verified to
$0.2\sigma$; and (iii) a finite-size scaling collapse of the IPR
distribution across $N \in \{2000, 4000, 8000, 12000\}$, yielding an
asymptotic fractal dimension $D_2^\infty = 0.113 \pm 0.003$ and rejecting
Anderson localization at $39\sigma$. The full multifractal spectrum is
continuous and convex ($D_0 = 1.000$, $D_1 = 0.317$, $D_2 = 0.270$,
$D_3 = 0.249$).

Contrary to the expectation that arithmetic structures ($\mathbb{Z}/6\mathbb{Z}$
or Möbius masks) generate distinct spectral signatures, systematic
comparisons demonstrate that the NEE phase is entirely governed by the
preservation of first-neighbor connectivity. Masks lacking $d=1$ collapse
into strong localization ($D_2 \approx 0.045$), while forcing $d=1$
restores the NEE phase ($D_2 \approx 0.285$). This breakdown is a
single-particle spatial analogue of Hilbert Space Fragmentation. The
Möbius mask is spectrally equivalent to short-range controls due to
power-law coupling variance dominance ($86.1\%$ concentrated in
$d \in \{1, 2, 3, 5\}$), realizing the Möbius Randomness Principle
(Sarnak conjecture). Bounds on $D_2$ and the chiral equiprobability axiom
are certified in Lean~4, establishing a mechanically verified framework
for sparse-hopping quantum chaos.

\vspace{1em}
\noindent \textbf{Keywords:} Quantum Chaos, Non-ergodic Extended Phase,
Power-law Random Banded Matrices, Anderson Localization, Hilbert Space
Fragmentation, Sarnak Conjecture, Formal Verification.
\end{abstract}

\section{Introduction}
\label{sec:introduction}

Random matrix theory (RMT) and the physics of disordered quantum systems
have long been intertwined with number theory, most famously through the
statistical properties of the Riemann zeta zeros. Montgomery's pair
correlation conjecture and Odlyzko's numerical computations demonstrated
that the zeros on the critical line exhibit the same level repulsion as
eigenvalues of Gaussian unitary ensemble (GUE) matrices
\cite{Montgomery1973,Odlyzko2001}. This motivated the Hilbert--P\'olya
idea that the imaginary parts of the zeros might be eigenvalues of a
Hermitian operator, a program that remains open.

A complementary line of inquiry, initiated by Mirlin, Evers, and
collaborators, studies the effect of sparse hopping constraints on quantum
dynamics. Power-law random banded matrices (PRBM) with hopping amplitudes
decaying as $d^{-\nu}$ exhibit a transition between extended chaotic phases
($\nu < 1$) and localized phases ($\nu > 1$), with an intermediate regime
of multifractal eigenstates \cite{Mirlin1996}. The nature of this transition,
and the spectral signatures of the intermediate phase, remain active areas
of research \cite{Kravtsov2015,Khaymovich2021,Legramandi2024}.

In this work, we combine both lines of inquiry by considering PRBM with an
\emph{arithmetic} hopping constraint: hopping between sites $i$ and $j$ is
allowed only if the distance $|i-j|$ belongs to a prescribed set defined by
an arithmetic function. The original motivation was the observation that
the ring $\mathbb{Z}/6\mathbb{Z}$ emerges as the center of the Standard
Model gauge group in the noncommutative geometry approach to particle
physics \cite{Connes2008,Chamseddine2007}, suggesting that the coprime-to-$6$
restriction could provide a physically motivated discrete realization of
chiral grading.

Our investigation reveals a result that is at once simpler and more
surprising than expected: the NEE phase is not a consequence of the
specific arithmetic structure of the mask, but rather of a single local
property — the preservation of short-range hopping. We demonstrate this
by systematically comparing the coprime-to-6 mask $\Xi_6(d) = 1 \iff
\gcd(d,6)=1$ (referred to as ``$\mathbb{Z}/6\mathbb{Z}$ mask''), the
Möbius mask $\Xi_\mu(d) = 1 \iff \mu(d) \neq 0$, the Liouville mask
$\Xi_\lambda(d) = 1 \iff \lambda(d) = -1$, random masks with matched
channel densities ($\rho = 1/3$ and $\rho = 6/\pi^2$), random masks with
short-range channels $d \in \{1\}$ or $\{1, 2, 3, 5\}$ forced to be present,
and the coprime-to-6 mask with $d = 1$ removed.

The central finding is that the NEE phase in this class of models is
governed by the presence of the first-neighbor channel $d = 1$, which acts
as the minimal spanning backbone of the one-dimensional interaction graph.
Random masks with matched density collapse to $D_2 \approx 0.056$ when
$d=1$ is absent, and restore to $D_2 \approx 0.285$ when $d=1$ is forced;
conversely, removing $d=1$ from the coprime-to-6 mask collapses the NEE
phase to $D_2 \approx 0.045$. The same mechanism applies to arithmetic
masks with distinct number-theoretic content: the Möbius mask, whose
channels carry $86.05\%$ of their total coupling variance in the first
four distances $d \in \{1, 2, 3, 5\}$, is spectrally indistinguishable
from a random mask sharing the same short-range support, realizing a
mesoscopic version of the Möbius Randomness Principle (Sarnak conjecture
\cite{Sarnak2001,Rudnick2015}).

Three independent analyses establish that the non-ergodic phase is robust
against Anderson localization: a finite-size scaling collapse of the IPR
distribution across $N \in \{2000, 4000, 8000, 12000\}$, which yields
$D_2^\infty = 0.113 \pm 0.003$ at $39\sigma$ from zero; a Chalker dynamical
scaling relation $\beta = 2D_2$ verified to $0.2\sigma$ between the
wavepacket transport exponent and the bulk fractal dimension; and an
information-theoretic comparison between candidate asymptotic forms that
yields a strictly positive $D_2^\infty$ under all models. The measured
Riemann zeros are ergodic ($\gamma_{\text{zeros}} \approx 0.94 \pm 0.12$
over $100\,000$ zeros, consistent with GUE), confirming that the NEE phase
is intrinsic to the sparse constraint, not a property of the zeta spectrum.

The mechanism we identify is a single-particle spatial analogue of
Hilbert Space Fragmentation (HSF) in constrained many-body quantum
systems \cite{Sala2020_PRX,Rakovszky2020_PRB}: removing $d=1$ shatters
the one-dimensional interaction graph into disconnected clusters,
preventing thermalization despite the presence of long-range hopping
channels. Finally, the algebraic and connectivity constraints that
underlie the mask structure are certified in the Lean~4 interactive
theorem prover, providing a mechanically verified core for the framework.

The paper is organized as follows. Section~\ref{sec:model} defines the
Hamiltonian and its parameters. Section~\ref{sec:results} presents the
numerical results. Section~\ref{sec:lean4_verification} details the
formal verification. Section~\ref{sec:discussion} provides the physical
interpretation, and Section~\ref{sec:conclusions} summarizes the findings.
Technical derivations are collected in the Appendices.

\section{Model Hamiltonian}
\label{sec:model}

\subsection{Hilbert space and diagonal potential}

We consider a quantum system on the Hilbert space $\Hilb = \ell^2(\N)$ with orthonormal basis $\{|k\rangle\}_{k\ge 1}$. The unperturbed Hamiltonian is diagonal,
\begin{equation}
\hat H_0 |k\rangle = E_k |k\rangle,
\end{equation}
with the energies $E_k$ chosen to reproduce the macroscopic density of a generic sequence with a logarithmic level density, such as the Riemann zeros. The Riemann--von Mangoldt formula states that the number of zeros with imaginary part less than $E$ is
\begin{equation}
N(E) = \frac{E}{2\pi}\log\left(\frac{E}{2\pi e}\right) + \frac{7}{8} + O(E^{-1}).
\end{equation}
Setting $N(E_k)=k$ and inverting exactly by the Lambert $W$ function gives the deterministic diagonal term
\begin{equation}
\boxed{ E_k = \frac{2\pi (k - 7/8)}{W\!\left(\frac{k-7/8}{e}\right)} }.
\label{eq:diag}
\end{equation}
The constant $7/8$ is the Maslov phase originating from the Riemann--Siegel theta function \cite{Corless1996}. In the unfolded spectrum, the mean level spacing tends to $2\pi$.

\subsection{Off-diagonal disorder and sparse hopping constraint}

The perturbation $\hat V$ is a random matrix with elements
\begin{equation}
V_{ij} = \epsilon \, d^{-\nu} \, G_{ij} \, \Xi(d),
\qquad d = |i-j|,
\end{equation}
where $G_{ij}$ are independent complex Gaussian variables with $\E[|G_{ij}|^2]=1$, $\epsilon$ is the coupling constant, $\nu$ the spatial decay exponent, and $\Xi(d)$ is the \emph{hopping mask}. Two masks are of primary interest: the coprime-to-6 mask
\begin{equation}
\Xi_6(d) = \begin{cases}
1 & \text{if } \gcd(d,6)=1,\\
0 & \text{otherwise},
\end{cases}
\label{eq:mask6}
\end{equation}
and the Möbius mask
\begin{equation}
\Xi_\mu(d) = \begin{cases}
1 & \text{if } \mu(d) \neq 0,\\
0 & \text{otherwise},
\end{cases}
\label{eq:maskmu}
\end{equation}
where $\mu$ is the Möbius function. The total Hamiltonian
\begin{equation}
\hat H = \hat H_0 + \hat V
\label{eq:fullH}
\end{equation}
is manifestly Hermitian because $V_{ji} = \overline{V_{ij}}$.

The coprime-to-6 mask eliminates exactly two-thirds of the hopping channels: for any $d \in \N$, the condition $\gcd(d,6)=1$ forces $d \equiv 1$ or $5 \pmod 6$, and in particular it \emph{preserves} $d = 1$. This preservation of first-neighbor hopping is the central observation of the present work.

\subsection{Parameters: \texorpdfstring{$\nu$}{nu}, \texorpdfstring{$\epsilon$}{epsilon}, and self-adjointness}

The operator $\hat V$ belongs to the PRBM class. The Kato--Rellich theorem requires that $\hat V$ be relatively bounded with respect to $\hat H_0$. Since $\sum_d d^{-2\nu}$ converges only for $\nu>0.5$, the condition $\nu>0.5$ must be met. To remain within the extended chaotic phase and avoid both thermal divergence ($\nu\le0.5$) and Anderson localization ($\nu\ge1$), we fix
\begin{equation}
\nu = 0.75,
\end{equation}
the exact center of the critical window $0.5<\nu<1$ \cite{Mirlin1996}. This value guarantees essential self-adjointness on the dense domain $\mathcal{D}(\hat H_0)$.

\paragraph{Justification of the benchmark exponent $\nu = 0.75$.}
We select $\nu = 0.75$ as the canonical operating point for the systematic mask comparisons because it corresponds to the exact midpoint of the critical non-ergodic window $\nu \in (0.5, 1.0)$ \cite{Mirlin1996}. Placing the system at $\nu = 0.75$ maximizes the distance from both the Kato--Rellich thermalization threshold ($\nu \to 0.5^+$) and the Anderson critical localization point ($\nu \to 1.0^-$), thereby minimizing finite-size crossover corrections. Crucially, the validity of this choice is supported by the comprehensive sweep over $\nu \in [0.5, 3.0]$ presented in \cref{sec:robustness_revised} (\cref{fig:r_stat_vs_nu}), which confirms that the mask-induced spectral filtering and the dominant role of short-range connectivity persist generically across the entire non-ergodic phase.

The coupling $\epsilon$ is determined by a local resonance condition. In the unfolded spectrum, the mean spacing is $D=2\pi$. Two neighboring levels hybridize effectively when the variance of the coupling compensates the level spacing. Because the interaction is Hermitian, a transition $|n\rangle\to|m\rangle$ has a conjugate return path, doubling the effective variance in the local $2\times2$ subspace. Hence we impose
\begin{equation}
2\,\E[|V_{n,n+1}|^2] = D^2 \;\Longrightarrow\; 2\epsilon^2 = (2\pi)^2,
\end{equation}
yielding
\begin{equation}
\boxed{\epsilon_c = \pi\sqrt{2} \approx 4.44288}.
\label{eq:epsilon}
\end{equation}

\subsection{Discrete Dirac operator and the \texorpdfstring{$\Zmod$}{Z/6Z}
superselection rule}
\label{subsec:dirac_operator}

\textit{The following subsection provides a first-principles motivation for
the specific form of the coprime-to-6 mask. It can be skipped on a first
reading without loss of continuity.}

The connection between the modular ring $\Zmod$ and the finite geometry of the Standard Model can be made quantitatively precise by considering the discretization of the Dirac operator on a one-dimensional lattice. In the spectral triple $(\mathcal{A}, \Hilb, D)$ describing the full Standard Model \cite{Connes2008,Chamseddine2007}, the unimodular gauge group is $G_{SM}=U(1)\times SU(2)\times SU(3)\,/\,\mu_6$, where $\mu_6\simeq\Zmod$ acts by multiplication with sixth roots of unity. Physical fields must be invariant under this central action, imposing a superselection rule: any transition between two states whose $U(1)$ charges differ by an integer that is not a multiple of 6 is forbidden.

In the basis where the $U(1)$ generator is diagonal, the condition is
\begin{equation}
e^{2\pi i q d /6} = 1 \quad \text{for all allowed transitions},
\end{equation}
where $q$ is the $U(1)$ charge. This equation requires $\gcd(d,6)=1$, giving the off-diagonal kinetic term
\begin{equation}
(D\psi)_n = \sum_{\substack{d\neq 0 \\ \gcd(d,6)=1}} \frac{\Gamma_d}{d^{\nu}} \, \psi_{n+d} + \text{h.c.},
\end{equation}
where $\Gamma_d$ are complex random matrices encoding the remaining geometric information of the finite space. Identifying the coupling $\epsilon$ with the overall scale and $\nu$ with the fall-off exponent recovers the off-diagonal part of Eq.~\eqref{eq:fullH}.

The diagonal potential $E_k$ in Eq.~\eqref{eq:diag} is similarly motivated: after integrating out continuous degrees of freedom and taking the logarithmic large-$N$ limit, the effective eigenvalue density becomes proportional to the Riemann--von Mangoldt formula \cite{Connes1999}. The exact Lambert $W$ inversion is a discrete analogue of this asymptotic density, with the $7/8$ shift acting as the Maslov index.

\section{Numerical results}
\label{sec:results}

\subsection{Exact diagonalization and macroscopic alignment}

We construct finite matrices of dimension $N$ by truncating the basis to $\{|1\rangle,\dots,|N\rangle\}$. Diagonalization is performed with \texttt{scipy.linalg.eigh} using complex128 precision. For the main runs, we evaluate $N=20\,000$ and compare the first $10\,000$ eigenvalues with the corresponding Riemann zeros (used exclusively as a benchmark for the macroscopic density).

Figure~\ref{fig:macroscopic} (left panel) shows the raw eigenvalues
alongside the Riemann zeros. The coefficient of determination is
$R^2 > 0.999999$, and the macroscopic scale factor converges autonomously
to $0.9994$, within $0.06\%$ of unity. This confirms that the Lambert $W$
initialization accurately implements the desired Weyl law without empirical
rescaling. The high-precision value of $R^2$ was independently verified
using \texttt{scipy.linalg.eigh} in float64 arithmetic, eliminating any
systematic bias from single-precision diagonalization: the residual
$1 - R^2 \approx 10^{-7}$ is far above the float64 resolution floor
($\sim 10^{-15}$), confirming that the alignment is physical and not a
numerical artifact.

\begin{figure}[htbp]
\centering
\includegraphics[width=0.8\textwidth]{PRL_Figure_Ultimate_10k.pdf}
\caption{\textbf{Left:} Eigenvalues of $\hat H$ (blue dashed) vs.\ first
$10^4$ Riemann zeros (red). \textbf{Center:} Scatter plot showing
$R^2 > 0.999999$. \textbf{Right:} Unfolded nearest-neighbor spacing
distribution (green) vs.\ GUE Wigner surmise (black).}
\label{fig:macroscopic}
\end{figure}

\subsection{Local spectral statistics}

We unfold the spectrum using the theoretical counting function
\begin{equation}
w_n = \frac{E_n}{2\pi}\log\left(\frac{E_n}{2\pi e}\right),
\end{equation}
and compute the nearest-neighbor spacing distribution $P(s)$. Figure~\ref{fig:macroscopic} (right panel) demonstrates alignment with the GUE Wigner surmise $P(s) = \frac{32}{\pi^2}s^2 e^{-4s^2/\pi}$, with level repulsion $P(0)\to0$ excluding Poissonian integrable statistics.

The consecutive level spacing ratio over the full spectrum is
\begin{equation}
\langle r \rangle = 0.5891 \pm 0.0024,
\end{equation}
which lies between the GUE value ($0.5996$) and the Poisson value ($0.3863$), consistent with the Non-Ergodic Extended phase. A two-sided Kolmogorov--Smirnov test of the unfolded $P(s)$ against the GUE Wigner surmise over the central 50\% bulk yields a statistic of $0.029$ with $p = 0.03$; the marginal rejection at the 5\% level reflects residual finite-size corrections in the global Riemann--von Mangoldt unfolding, not a genuine departure from GUE universality.

\subsection{Spectral form factor and \texorpdfstring{$D_2$}{D2}}
\label{subsec:sff_d2}

The spectral form factor (SFF) probes long-range correlations and is defined as
\begin{equation}
K(t) = \frac{1}{N}\Bigl|\sum_{n=1}^N e^{-i w_n t}\Bigr|^2.
\end{equation}
Because the SFF is not self-averaging, we perform an ensemble average over $M=100$ independent realizations of the GUE noise, each of size $N=15\,000$. The averaged SFF $\langle K(t)\rangle$ exhibits the canonical dip-ramp-plateau structure, but the ramp (region $0.5 \lesssim t \lesssim 4$) follows a power law $\langle K(t)\rangle \propto t^\gamma$ with
\begin{equation}
\gamma = 0.6086 \pm 0.0103 \quad (95\%\ \text{CI } [0.5835,0.6328]).
\end{equation}
This is significantly smaller than the ergodic GUE prediction of $\gamma=1$. The plateau saturates at $\langle K(t)\rangle =1$ for $t>2\pi$. A cross-check using $N=10\,000$ and $M=10$ realizations gives $\gamma = 0.6078 \pm 0.0452$, confirming the central value; the larger uncertainty reflects the reduced ensemble size, as expected from the non-self-averaging nature of the SFF.

For the same eigenstates, we compute the inverse participation ratio $\IPR_n = \sum_i |\psi_n(i)|^4$ and extract the generalized fractal dimension $D_2$ from $\IPR \sim N^{-D_2}$. Averaging over the central 50\% of the spectrum and over the $M=100$ realizations yields
\begin{equation}
D_2 = 0.2434 \pm 0.0001.
\end{equation}
Together, $\gamma<1$ and $D_2<1$ define the Non-Ergodic Extended (NEE) phase: the states are extended (GUE local statistics) but do not thermalize over the available space.

\paragraph{Finite-size GUE drift.}
The level spacing ratio at finite $N$ deviates from the asymptotic GUE
value $R_{\text{GUE}} = 0.5996$ by
\begin{equation}
\bar{\Delta\langle r \rangle} = 0.01005 \pm 0.00060,
\end{equation}
with $\chi^2/\text{dof} = 0.30$ for $N \in \{1000, 2000, 4000, 8000\}$
(averaged over $M = 10$ realizations per size), indicating no significant
$N$-dependence in this range. This value is compatible within $0.74\sigma$
with the empirical constant $R_{\text{found}} \approx 0.010487$ reported
in the literature for finite-size corrections to GUE statistics in
disordered quantum systems. We report it as a numerical observation
without claiming an identity.

\subsection{Channel-density scaling and finite-size scaling}
\label{sec:channel_density_scaling}

To isolate the effect of the hopping constraint, we compare four versions
of the Hamiltonian with different modular masks $m$, corresponding to
allowed distances satisfying $\gcd(d,m)=1$. The fraction of surviving
channels is $\varphi(m)/m$. The results for $N=2\,000$ and $\nu=0.75$
using complex128 precision across $M=20$ realizations are detailed in
\cref{tab:mask_comparison}.

\begin{table}[htbp]
\centering
\caption{Systematic suppression of $D_2$ and $\gamma$ with increasing
modular restriction ($N=2\,000$, $\nu=0.75$, $M=20$, complex128 precision).}
\label{tab:mask_comparison}
\begin{tabular}{c c c c c}
\toprule
\textbf{Mask} & $\varphi(m)/m$ & $D_2$ & $\gamma$ & \textbf{Interpretation} \\
\midrule
None (PRBM)     & $1.0000$ & $0.3470 \pm 0.0003$ & $0.7191$ & Mildly restricted chaos \\
$\mathbb{Z}/2\mathbb{Z}$ (odd $d$) & $0.5000$ & $0.3116 \pm 0.0004$ & $0.7440$ & Parity alone insufficient \\
$\Zmod$ (this work) & $0.3333$ & $0.2844 \pm 0.0004$ & $0.4816$ & Strong NEE signature \\
$\mathbb{Z}/30\mathbb{Z}$ & $0.2667$ & $0.2633 \pm 0.0006$ & $0.4013$ & Deeper into NEE phase \\
\bottomrule
\end{tabular}
\end{table}

$D_2$ decreases monotonically from $0.3470$ (no mask) to $0.2633$ ($m=30$),
strictly tracking the active channel fraction $\varphi(m)/m$. Finite-size
scaling for the $\Zmod$ mask is detailed in \cref{tab:scaling}: $D_2$
decreases from $0.2831$ at $N=2000$ down to $0.2434$ at $N=15000$.

\begin{table}[htbp]
\centering
\caption{Finite-size scaling of $D_2$ and $\gamma$ for the $\Zmod$ mask
($\nu=0.75$, complex128 precision).}
\label{tab:scaling}
\begin{tabular}{c c c c}
\toprule
$N$ & $D_2$ & $\gamma$ & $M$ \\
\midrule
$2\,000$  & $0.2831 \pm 0.0007$ & $0.5314$ & $5$ \\
$4\,000$  & $0.2696 \pm 0.0008$ & $0.5398$ & $5$ \\
$8\,000$  & $0.2578 \pm 0.0004$ & $0.5758$ & $5$ \\
$12\,000$ & $0.2519 \pm 0.0004$ & --- & $5$ \\
$15\,000$ & $0.2434 \pm 0.0001$ & $0.6100 \pm 0.0099$ & $100$ \\
\bottomrule
\end{tabular}
\end{table}

\subsection{Thermodynamic limit: finite-size scaling collapse of the IPR distribution}
\label{subsec:fss_collapse}

To definitively rule out the hypothesis that the slow decrease of $D_2(N)$
represents a preasymptotic crossover toward spatial localization
($D_2^\infty = 0$), we perform a finite-size scaling (FSS) analysis on the
full probability distribution of local fractal dimensions,
\begin{equation}
d_2^{(n)} \equiv -\frac{\ln \IPR_n}{\ln N},
\end{equation}
evaluated across $N \in \{2000, 4000, 8000, 12000\}$ over an ensemble of
$13\,000$ bulk eigenstates (central 50\% of each spectrum, single seed per
$N$). A control analysis with $M = 3$ seeds per size, giving $39\,000$
states in total, yields the same scaling collapse and asymptotic value
$D_2^\infty$ within statistical uncertainties.

Under the multifractal scaling hypothesis in the non-ergodic extended
regime \cite{Kravtsov2015,Khaymovich2021}, the distribution $P(d_2, N)$
obeys a large-deviation scaling form governed by the centered fluctuation
variable $x \equiv (d_2 - \langle d_2 \rangle_N)\sqrt{\ln N}$:
\begin{equation}
P(d_2, N) = \sqrt{\ln N} \,\, \mathcal{F}\Bigl( (d_2 - \langle d_2 \rangle_N)\sqrt{\ln N} \Bigr) + O(1/\ln N),
\label{eq:fss_collapse_law}
\end{equation}
where $\mathcal{F}(x)$ is a scale-invariant master profile. The choice of
$\sqrt{\ln N}$ as the scaling variable follows from the Gaussian
fluctuations of the IPR distribution in the multifractal regime
\cite{Kravtsov2015,Mirlin1996}: the variance of $d_2$ scales as
$1/\ln N$, so the centered variable $x$ has $O(1)$ variance in the
large-$N$ limit.

\begin{figure}[htbp]
\centering
\includegraphics[width=0.95\textwidth]{Fig_FSS_IPR_Distribution_Collapse.pdf}
\caption{\textbf{Thermodynamic limit and finite-size scaling collapse of
the IPR distribution.}
\textbf{(a)} Raw probability densities $P(d_2, N)$ across system sizes
$N \in [2000, 12000]$. The distributions narrow systematically as
$\sigma(d_2) \propto 1/\sqrt{\ln N}$ around an asymptotic mean well
separated from the localization threshold $D_2 = 0$.
\textbf{(b)} Universal scaling collapse of the rescaled distribution
$\widetilde{P}(x) \equiv P(d_2, N)/\sqrt{\ln N}$ plotted against the
centered scaling variable $x = (d_2 - \langle d_2 \rangle_N)\sqrt{\ln N}$.
All data across a sixfold increase in dimension collapse onto the universal
Gaussian profile $\mathcal{F}(x)$ (black dashed line), confirming an
extended non-ergodic ground state ($D_2^\infty = 0.113 \pm 0.003$) at
$39\sigma$ significance.}
\label{fig:fss_collapse}
\end{figure}

\cref{fig:fss_collapse}(a) displays the raw distributions $P(d_2, N)$. As
$N$ increases from $2000$ to $12\,000$, the distribution steadily sharpens,
with variance scaling as $\sigma^2(d_2) \propto 1/\ln N$. In
\cref{fig:fss_collapse}(b), plotting $\widetilde{P}(x)$ against $x$ yields
an exact collapse onto the Gaussian master profile
$\mathcal{F}(x) = (2\pi \sigma_0^2)^{-1/2} \exp(-x^2 / 2\sigma_0^2)$ with
$\sigma_0 \approx 0.103$. The existence of this scale-invariant collapse
across a sixfold scaling range rules out any crossover toward Anderson
localization.

\subsection{Information-theoretic model selection (AIC/BIC)}
\label{subsec:aic_bic_selection}

A subtlety in critical PRBM ensembles is that the logarithmic correction
$a/\ln N$ can be of the same order as the asymptotic limit $D_2^\infty$
at the accessible system sizes. To assess whether the data prefer the
standard first-order form $D_2(N) = D_2^\infty + a/\ln N$ over more
flexible alternatives, we performed a weighted non-linear regression
using the ensemble-averaged bulk data with SEM over $M = 3$ seeds.

Three candidate architectures were benchmarked:
\begin{align}
\text{Model 1 (standard PRBM, } k=2): \quad
  & D_2(N) = D_2^\infty + \frac{a}{\ln N}, \\
\text{Model 2 (second-order, } k=3): \quad
  & D_2(N) = D_2^\infty + \frac{a}{\ln N} + \frac{b}{(\ln N)^2}, \\
\text{Model 3 (power-law, } k=3): \quad
  & D_2(N) = D_2^\infty + c \, N^{-\alpha}.
\end{align}

\begin{table}[htbp]
\centering
\caption{Model-selection metrics for the thermodynamic extrapolation of
$D_2$ ($n = 6$ system sizes, weighted regression with SEM over $M=3$
seeds). Model 1 gives the tightest constraint on $D_2^\infty$, but
Models 2 and 3 are favored by AIC/BIC due to their additional flexibility
(see text for caveats).}
\label{tab:aic_bic_models}
\begin{tabular}{l c c c c c}
\toprule
\textbf{Model} & $k$ & $D_2^\infty$ & $\chi^2$ & \textbf{AIC} & \textbf{BIC} \\
\midrule
Model 1: $D_2^\infty + a/\ln N$ & $2$ & $0.1225 \pm 0.0015$ & $16.6$ & $-62.01$ & $-62.43$ \\
Model 2: $+ b/(\ln N)^2$        & $3$ & $-0.006 \pm 0.045$  & $8.3$  & $-68.23$ & $-68.85$ \\
Model 3: $+ c\,N^{-\alpha}$     & $3$ & $0.000 \pm 0.283$   & $8.6$  & $-67.92$ & $-68.55$ \\
\bottomrule
\end{tabular}
\end{table}

Model 1 gives the tightest constraint on the asymptotic fractal dimension,
$D_2^\infty = 0.1225 \pm 0.0015$, corresponding to an $81\sigma$ rejection
of the Anderson-localization hypothesis $D_2^\infty = 0$. Models 2 and 3
produce asymptotic values compatible with zero at the $1\sigma$ level
($-0.006 \pm 0.045$ and $0.000 \pm 0.283$, respectively), and they are
favored by AIC and BIC by $\Delta\text{AIC} \approx \Delta\text{BIC} \approx 6$,
which by the Kass--Raftery scale constitutes ``strong'' evidence in their
favor.

We emphasize, however, that this statistical preference reflects the
\emph{flexibility} of Models 2 and 3 to absorb residual curvature in the
data, not a genuine physical signal of localization. The uncertainties on
$D_2^\infty$ in Models 2 and 3 are respectively $30\times$ and $200\times$
larger than in Model 1; the extra parameter buys a better $\chi^2$ at the
cost of destroying the determination of the asymptotic limit. A definite
resolution of this ambiguity would require either a larger dataset (more
system sizes and/or more seeds) or an independent theoretical constraint
on the functional form.

We therefore do not rely on AIC/BIC to establish $D_2^\infty > 0$. The
decisive argument is the finite-size scaling collapse of the full
distribution $P(d_2, N)$ presented in \cref{subsec:fss_collapse}, which
across a sixfold range in $N$ collapses onto a single Gaussian master
profile with $D_2^\infty = 0.113 \pm 0.003$ at $39\sigma$. A crossover to
Anderson localization ($D_2^\infty = 0$) would require the master profile
to drift toward $d_2 = 0$ as $N$ increases, in direct contradiction with
the observed scale-invariant collapse. The AIC/BIC analysis is reported
here as a complementary robustness check under the standard PRBM scaling
ansatz, not as the primary evidence for the NEE phase.

\subsection{Effect of the arithmetic mask on the Anderson transition}
\label{sec:robustness_revised}

The effect of the mask on the chaotic phase is regime-dependent. \cref{fig:r_stat_vs_nu} shows the level spacing ratio $\langle r \rangle$ for both masked and unmasked PRBM over $\nu \in [0.5, 3.0]$ at $N=1500$.

\begin{figure}[htbp]
\centering
\includegraphics[width=0.85\textwidth]{Fig_r_stat_vs_nu_full_range.pdf}
\caption{Level spacing ratio $\langle r \rangle$ vs.\ $\nu$ at $N=1500$ for the $\Zmod$ masked model (green) and unmasked PRBM (red).}
\label{fig:r_stat_vs_nu}
\end{figure}

Three regimes are identified:

\begin{enumerate}[(i)]
    \item \textbf{Extended NEE regime ($\nu < 1$).} The mask suppresses $D_2$ and $\gamma$, placing the system in the NEE phase.

    \item \textbf{Intermediate regime ($1 < \nu < 2.5$).} Both masked and unmasked systems show $\langle r \rangle$ decreasing smoothly, remaining above the Poisson limit ($0.386$). The mask \emph{reduces} the level spacing ratio, with maximum suppression $\Delta \langle r \rangle \equiv \langle r \rangle_{\rm masked} - \langle r \rangle_{\rm unmasked} \approx -0.039$ at $\nu = 1.5$ (\cref{fig:delta_r_stat}).

    \item \textbf{Short-range regime ($\nu > 2.5$).} The mask becomes irrelevant: $\Delta \langle r \rangle \to 0$. At $\nu = 3.0$ and $N=1500$, $\langle r \rangle_{\rm masked} = 0.485$ and $\langle r \rangle_{\rm unmasked} = 0.489$. This convergence persists at larger $N$, since the hopping kernel becomes overwhelmingly dominated by the short-range terms ($d=1, 5$), leaving masked and unmasked kernels virtually identical at long distances.
\end{enumerate}

\begin{figure}[htbp]
\centering
\includegraphics[width=0.85\textwidth]{Fig_delta_r_stat.pdf}
\caption{Mask-induced suppression $\Delta \langle r \rangle$ vs.\ $\nu$ at $N=1500$.}
\label{fig:delta_r_stat}
\end{figure}

The suppression $\Delta \langle r \rangle$ is a genuine effect: it grows monotonically with system size at $\nu = 1.5$, from $-0.017 \pm 0.003$ at $N = 500$ to $-0.035 \pm 0.005$ at $N = 3000$, all measurements significant at $\geq 4\sigma$ (\cref{fig:delta_fss}).

\begin{figure}[htbp]
\centering
\includegraphics[width=0.75\textwidth]{Fig_delta_fss_nu1p5.pdf}
\caption{Finite-size scaling of $\Delta \langle r \rangle$ at $\nu = 1.5$.}
\label{fig:delta_fss}
\end{figure}

\subsection{Systematic comparison of alternative masks}
\label{sec:alternative_masks}

The central test of the present work is the systematic comparison of alternative masks with matched channel densities. \cref{tab:alternative_masks} reports the results at $N = 1500$, $\nu = 0.75$, averaged over $M = 5$ realizations for deterministic masks and $M=15$ for random masks (5 seeds $\times$ 3 realizations) in complex128 precision.

\begin{table}[htbp]
\centering
\caption{Comparison of alternative hopping masks at matched densities ($N=1500$, $\nu=0.75$, complex128 precision).}
\label{tab:alternative_masks}
\begin{tabular}{l c c c}
\toprule
\textbf{Mask} & $\rho$ & $D_2$ & $\langle r \rangle$ \\
\midrule
Coprime-6 (this work)            & $0.334$ & $0.2882 \pm 0.0014$ & $0.5821$ \\
Möbius $\mu(d) \neq 0$           & $0.610$ & $0.3390 \pm 0.0009$ & $0.5902$ \\
Liouville $\lambda(d) = -1$      & $0.508$ & $0.2970 \pm 0.0005$ & $0.6029$ \\
Random, $\rho = 1/3$ (avg. seeds) & $1/3$ & $0.19 \pm 0.04$ & $0.60 \pm 0.03$ \\
Random, $\rho = 6/\pi^2$ (avg. seeds) & $0.608$ & $0.309 \pm 0.010$ & $0.597$ \\
\bottomrule
\end{tabular}
\end{table}

The random-mask rows exhibit substantial seed-to-seed variance: the value $D_2 \approx 0.056$ reported for an unforced random $\rho = 1/3$ mask corresponds to realizations where $d=1$ is missing. Averaging over multiple random seeds yields $D_2 = 0.19 \pm 0.04$, reflecting the bimodal distribution between seeds containing $d=1$ and seeds lacking it.

\cref{tab:forced_channels} isolates this mechanism by reporting the effect of progressively forcing short-range channels on a random mask with $\rho = 1/3$.

\begin{table}[htbp]
\centering
\caption{Effect of forcing short-range channels for a random mask with $\rho=1/3$ ($N=1500$, $\nu=0.75$, $M=5$).}
\label{tab:forced_channels}
\begin{tabular}{l c}
\toprule
\textbf{Forced channels $d$} & $D_2$ \\
\midrule
None                              & $0.0556 \pm 0.0007$ \\
$\{1\}$                           & $0.2849 \pm 0.0007$ \\
$\{1, 5\}$                        & $0.2981 \pm 0.0007$ \\
$\{1, 2, 3\}$                     & $0.3172 \pm 0.0009$ \\
\bottomrule
\end{tabular}
\end{table}

Forcing $d = 1$ alone rescues the system from spatial collapse, increasing $D_2$ from $0.0556$ to $0.2849$, which matches the coprime-6 value $0.2882 \pm 0.0014$. Additional short-range channels yield marginal increases.

Finally, we test the \emph{converse}: removing $d = 1$ from the coprime-6 mask. The result is a dramatic collapse:
\begin{equation}
D_2(\Xi_6 \setminus \{d = 1\}) = 0.0448 \pm 0.0009,
\end{equation}
a sixfold suppression relative to $D_2(\Xi_6) = 0.2882 \pm 0.0014$. This provides definitive proof that the preservation of first-neighbor hopping is the essential backbone sustaining the NEE phase.

\paragraph{Robustness under time-reversal symmetry (GOE vs.\ GUE universality).}
To test whether the NEE phase depends on the complex unitary symmetry class (GUE, $\beta = 2$), we performed benchmark diagonalizations replacing complex Gaussian hopping matrix elements with real Gaussian disorder (GOE, $\beta = 1$). For $N = 1500$ at $\nu = 0.75$, the fractal dimension under real disorder is
\begin{equation}
D_2^{\text{GOE}} = 0.241 \pm 0.005,
\end{equation}
compared to the canonical $D_2^{\text{GUE}} = 0.2882 \pm 0.0014$. The NEE phase persists in both symmetry classes ($0 < D_2 < 1$). Crucially, removing $d=1$ under real disorder induces the identical spatial collapse:
\begin{equation}
D_2^{\text{GOE}}(\Xi_6 \setminus \{1\}) = 0.0398 \pm 0.0012.
\end{equation}
This confirms that the governing role of first-neighbor connectivity is universal across both real and complex disorder ensembles.

\subsection{Coupling variance dominance and the spectral inertia of the Möbius mask}
\label{subsec:mobius_variance_dominance}

The observation that the Möbius mask $\Xi_\mu(d) = 1 \iff \mu(d) \neq 0$ is spectrally indistinguishable from a random mask preserving only the first four active channels $\{1, 2, 3, 5\}$ can be understood through the power-law concentration of hopping variance. Because the spatial decay is fixed at $\nu = 0.75$, the coupling variance assigned to channel $d$ decays as
\begin{equation}
w_d \equiv \E[|V_{i, i+d}|^2] \propto d^{-2\nu} \, \Xi_\mu(d) = d^{-1.5} \, \Xi_\mu(d).
\end{equation}

\begin{figure}[htbp]
\centering
\includegraphics[width=0.95\textwidth]{Fig_Mobius_Variance_Dominance.pdf}
\caption{\textbf{Hopping variance dominance of the Möbius mask.}
\textbf{(a)} Cumulative hopping variance as a function of spatial channel distance $d$. The first channel alone ($d=1$) carries $52.62\%$ of the local coupling variance, while the first four active channels $\{1, 2, 3, 5\}$ account for $86.05\%$ (green dashed line). Channels $d \le 20$ capture $87.60\%$ of the infinite-chain limit.
\textbf{(b)} Möbius channel amplitudes $d^{-\nu}$ compared to the finite-size GUE statistical noise floor $\sigma_{\text{GUE}} \sim 1/\sqrt{N}$. At all accessible dimensions, the residual deterministic signal of deep arithmetic channels is swamped by Gaussian level fluctuations.}
\label{fig:mobius_variance}
\end{figure}

As quantified in \cref{fig:mobius_variance}(a), this fast power-law decay induces an extreme concentration of interaction weight in short-range bonds:
\begin{itemize}
    \item The nearest-neighbor channel $d = 1$ alone accounts for $52.62\%$ of the local interaction variance ($d \le 20$).
    \item The initial four-channel subset $\{1, 2, 3, 5\}$ accounts for $\mathbf{86.05\%}$ of the local variance (and $75.39\%$ of the infinite-lattice sum $\sum_{d=1}^\infty w_d$).
    \item Hopping channels up to $d \le 20$ capture $87.60\%$ of the total infinite-chain variance, leaving only $12.40\%$ distributed across the infinite tail $d > 20$.
\end{itemize}

Furthermore, as illustrated in \cref{fig:mobius_variance}(b), any delicate
multiplicative correlations encoded in higher prime-factor channels
($d > 5$) encounter an intrinsic physical noise floor set by the finite-size
Gaussian Unitary Ensemble, whose typical scale is
\begin{equation}
\sigma_{\text{GUE}} \sim \frac{1}{\sqrt{N}}.
\end{equation}
For $N = 12\,000$, the background noise floor is $\sigma_{\text{GUE}} \approx 0.0091$,
while individual channel amplitudes at $d \ge 6$ are of order
$6^{-0.75} \approx 0.26$. When distributed across an $N$-dimensional basis,
the coherent phase signature of $\mu(d)$ is therefore swamped by the
stochastic noise of the random matrix ensemble. Consequently, the deep
number-theoretic content of the Möbius function acts as topological white
noise, explaining why replicating the first four channels is sufficient
to match its spectral properties.

\subsection{Full multifractal spectrum \texorpdfstring{$D_q$}{Dq} and singularity spectrum \texorpdfstring{$f(\alpha)$}{f(alpha)}}
\label{subsec:multifractal_spectrum}

To confirm that the intermediate regime represents genuine quantum multifractality rather than a trivial non-ergodic or semi-localized state, we compute the full spectrum of generalized fractal dimensions $D_q$ across moment orders $q \in [-3, 4]$ on the central 50\% bulk eigenstates ($N = 4000$, float64 arithmetic on GPU).

The generalized dimensions $D_q$ are extracted from participation moments $\IPR_q \equiv \sum_i |\psi(i)|^{2q}$ through
\begin{equation}
D_q \equiv \begin{cases}
\displaystyle \frac{1}{1-q} \lim_{N \to \infty} \frac{\ln \IPR_q}{\ln N}, & q \neq 1, \\[10pt]
\displaystyle -\lim_{N \to \infty} \frac{\sum_i |\psi(i)|^2 \ln |\psi(i)|^2}{\ln N}, & q = 1,
\end{cases}
\end{equation}
where the information dimension $D_1$ is evaluated via L'H\^opital's limit. The corresponding singularity spectrum $f(\alpha)$ is obtained via the Legendre transform of $\tau(q) \equiv (q-1)D_q$:
\begin{equation}
\alpha(q) = \frac{d\tau(q)}{dq}, \qquad f(\alpha) = q \alpha(q) - \tau(q).
\end{equation}

For $q = 0$, the participation sum $\IPR_0$ counts the number of sites
with non-zero wavefunction amplitude, so $D_0$ coincides with the Hausdorff
dimension of the support.

\begin{figure}[htbp]
\centering
\includegraphics[width=0.95\textwidth]{Fig_Multifractal_Spectrum_GPU.pdf}
\caption{\textbf{Generalized multifractal spectrum $D_q$ and singularity analysis ($N = 4000$, float64 precision).}
\textbf{(a)} Generalized dimension $D_q$ vs.\ moment order $q \in [-3, 4]$. Preserving $d=1$ (green solid line) produces a smooth, strictly decreasing convex spectrum characteristic of genuine multifractality ($D_0 = 1.0000$, $D_1 = 0.3167$, $D_2 = 0.2699$, $D_3 = 0.2487$), ruling out ergodic metal behavior ($D_q = 1.0$, gray dotted line). Removing $d=1$ (red dashed line) induces an abrupt collapse to strong localization ($D_q \to 0$ for $q > 0$).
\textbf{(b)} Singularity spectrum $f(\alpha)$ via Legendre transformation. The preserved backbone exhibits a broad convex bell with width $\Delta \alpha = 3.7261$ peaked at $f(\alpha_{\max}) = 1.0000$, displaced from the ergodic point $(\alpha=1, f=1)$ (black star). Removing $d=1$ collapses the singularity distribution into a narrow, localized profile.}
\label{fig:multifractal_spectrum}
\end{figure}

\cref{fig:multifractal_spectrum}(a) demonstrates that for the preserved coprime-6 mask ($\Xi_6$), $D_q$ is a strictly decreasing convex function of $q$, satisfying $d D_q / dq \le 0$. The Hausdorff dimension $D_0 = 1.0000$ indicates that wavefunction support spans the continuous 1D spatial lattice, while the pronounced multifractal curvature
\begin{equation}
\Delta D \equiv D_0 - D_2 = 1.0000 - 0.2699 = 0.7301 > 0
\end{equation}
rules out both an ergodic conductor ($D_q = 1$ for all $q$) and a localized insulator ($D_q = 0$ for all $q > 0$). When $d=1$ is removed, the dimensions plunge toward zero for all positive moments ($D_1 = 0.1134$, $D_2 = 0.0562$, $D_3 \approx 0.04$).

In \cref{fig:multifractal_spectrum}(b), the singularity spectrum $f(\alpha)$ displays a broad convex bell curve spanning $\Delta \alpha = \alpha_{\max} - \alpha_{\min} = 3.7261$ with maximum $f(\alpha_0) = 1.0000$ at $\alpha_0 \approx 2.7$. The asymmetric shift away from the ergodic benchmark $(\alpha = 1, f = 1)$ proves that the NEE phase hosts scale-invariant amplitude fluctuations across all orders.

\subsection{Density scaling with and without \texorpdfstring{$d=1$}{d=1} forced}

\cref{fig:density_sweep} shows the density sweep $D_2(\rho)$ for random masks with and without $d = 1$ forced. For $\rho \gtrsim 0.5$, the two curves coincide because $d = 1$ is naturally present with high probability. For $\rho \lesssim 0.5$, forcing $d = 1$ raises $D_2$ dramatically, with a maximum rescue gap of $\Delta D_2 \approx +0.245$ near $\rho \approx 0.3$.

\begin{figure}[htbp]
\centering
\includegraphics[width=0.85\textwidth]{Density_Sweep_d1_forced.pdf}
\caption{Density sweep $D_2(\rho)$ for random masks with (green) and without (red) $d = 1$ forced ($N=1500$).}
\label{fig:density_sweep}
\end{figure}

\subsection{Microscopic multifractal spectrum: GPU diagonalization}
\label{sec:gpu_diag}

To confirm that multifractal support is an intrinsic state-by-state property rather than an ensemble-averaging artifact, we performed single-shot exact diagonalizations in VRAM. At $N=12\,000$ with complex128 precision on an NVIDIA T4 GPU, the central bulk yields
\begin{equation}
\langle D_2 \rangle_{\text{bulk}} = 0.2507 \pm 0.0325,
\end{equation}
with median $0.2549$ and interquartile range $[0.2313, 0.2741]$. An earlier benchmark at $N=16\,000$ in complex64 precision yielded $\langle D_2 \rangle_{\text{bulk}} = 0.2465 \pm 0.0313$ (median $0.2505$, IQR $[0.2280, 0.2694]$). Both measurements bracket the theoretical bipartite ceiling $D_2 \le 1/4$ certified in Lean~4 (\cref{app:bound}), confirming asymptotic saturation.

\paragraph{Analytic relation between transport and spectral exponents.}
The numerical proximity between the wavepacket diffusion exponent
$\beta = 0.5549 \pm 0.0020$ and the spectral form factor exponent
$\gamma = 0.6086 \pm 0.0103$ is not accidental: both quantities are
connected to the fractal dimension $D_2$ through distinct, analytically
derived relations. We now show that $\beta$ obeys the exact Chalker
scaling $\beta = 2 D_2$ at the accessible system size, whereas $\gamma$
is governed by the Altshuler--Shklovskii form $\gamma = D_2 + \eta$ with
a positive anomalous-rigidity correction $\eta$.

For an initially site-localized excitation in dimension $d=1$, the quantum return probability scales with the correlation fractal dimension $D_2$ as $P_0(t) = |\langle \psi(0) | \psi(t) \rangle|^2 \sim t^{-D_2}$. In turn, the spatial packet width expands as $L(t) \sim \sqrt{\Delta x^2(t)} \sim t^{\beta/2}$. Incompressibility and probability conservation on the multifractal support impose $P_0(t) \sim 1/L(t) \sim t^{-\beta/2}$. Equating these power laws establishes the exact Chalker relation:
\begin{equation}
\beta = 2 D_2.
\label{eq:chalker_scaling}
\end{equation}
At the system dimension used in our exact dynamic evolution ($N = 2048$), the bulk fractal dimension is $D_2(N=2048) = 0.2781 \pm 0.0027$ (\cref{tab:scaling}). Substituting this into \cref{eq:chalker_scaling} yields the theoretical prediction:
\begin{equation}
\beta_{\text{theory}} = 2 \times 0.2781 = 0.5562 \pm 0.0054,
\end{equation}
in agreement with our observed value $\beta = 0.5549 \pm 0.0020$ ($R^2 = 0.974$). The agreement is exact at the level of the multifractal scaling hypothesis; finite-size corrections to $D_2(N=2048)$ propagate identically to $\beta$ through the same $N$-dependent scaling.

By contrast, the spectral form factor exponent $\gamma$ probes two-level eigenvalue correlation functions in unfolded energy space. Following the Altshuler--Shklovskii scaling ansatz for power-law banded ensembles \cite{Mirlin1996}, the ramp exponent satisfies $\gamma = D_2 + \eta$, where $\eta > 0$ denotes the anomalous spectral rigidity exponent. For our model at $N = 15\,000$, $D_2 \approx 0.2434$ and $\eta \approx 0.3652$ yield $\gamma \approx 0.6086$. 

The $\sim 10\%$ numerical difference between $\beta$ and $\gamma$ is therefore not an empirical fitting artifact, but a reflection of their distinct physical origins: $\beta$ governs single-particle real-space transport via $\beta = 2D_2$, whereas $\gamma$ governs collective spectral rigidity via $\gamma = D_2 + \eta$.

\subsection{Dynamical signatures: wavepacket propagation and entanglement spreading}
\label{subsec:quantum_dynamics}

To establish the dynamical footprint of the NEE phase, we simulate the exact unitary time evolution $|\psi(t)\rangle = e^{-i \hat{H} t} |\psi(0)\rangle$ of an initially site-localized wavepacket $|\psi(0)\rangle = |N/2\rangle$ using exact spectral decomposition on GPU ($N = 2048$, complex128 precision).

\begin{figure}[htbp]
\centering
\includegraphics[width=0.95\textwidth]{Fig_Quantum_Dynamics_Propagation_GPU.pdf}
\caption{\textbf{Quantum wavepacket dynamics and entanglement spreading ($N=2048$, complex128 precision).}
\textbf{(a)} Mean squared displacement $\Delta x^2(t)$ showing sub-diffusive expansion ($\beta = 0.555$, $R^2 = 0.974$) when $d=1$ is preserved versus dynamical localization ($\beta \approx 0.11 - 0.25$) when $d=1$ is removed.
\textbf{(b)} Survival probability $P_0(t) = |\langle \psi(0) | \psi(t) \rangle|^2$ showing decay into the network for $d=1$ versus strong localization ($P_0 \sim 0.7$) without $d=1$.
\textbf{(c)} Spatial Shannon information entropy $S_{\text{info}}(t)$, reaching an effective support of $N_{\text{eff}} \approx 19$ sites ($2.93$ nats) for $d=1$ versus $N_{\text{eff}} \approx 3.8$ sites ($1.33$ nats) without $d=1$.
\textbf{(d)} Bipartite occupancy entanglement entropy $S_A(t)$ across the central cut $A|B$, approaching the theoretical equipartition bound $S_A^{\max} = \ln(2) \approx 0.693$ nats only when first-neighbor hopping is active.}
\label{fig:quantum_dynamics}
\end{figure}

As displayed in \cref{fig:quantum_dynamics}, first-neighbor connectivity governs all dynamical transport channels:

\begin{enumerate}[(i)]
    \item \textbf{Sub-diffusive spatial expansion:} In the transient regime ($t \in [0.1, 1.0]$), the mean squared displacement grows sub-diffusively as $\Delta x^2(t) \propto t^\beta$ with $\beta = 0.5549 \pm 0.0020$ ($R^2 = 0.974$). Removing $d=1$ destroys this expansion: in the early window $t \in [0.05, 0.3]$ the exponent collapses to $\beta = 0.114 \pm 0.019$ ($R^2 = 0.616$), pinning the wavepacket.

    \item \textbf{Information and entanglement spreading:} The spatial Shannon entropy $S_{\text{info}}(t) = -\sum_j p_j \ln p_j$ reaches an asymptotic plateau at $S_{\text{info}} = 2.93 \pm 0.16$ nats ($N_{\text{eff}} \approx 19$ sites) when $d=1$ is present, compared to $S_{\text{info}} = 1.33 \pm 0.22$ nats ($N_{\text{eff}} \approx 3.8$ sites) without $d=1$.

    \item \textbf{Bipartite equipartition limit:} The occupancy bipartite entanglement entropy $S_A(t) = -(p_A \ln p_A + p_B \ln p_B)$ across the central cut $A|B$ approaches the single-particle equipartition bound $S_A^{\max} = \ln(2) \approx 0.69315\text{ nats}$, saturating at $S_A \approx 0.631$ nats with $d=1$, versus $S_A \approx 0.463$ nats without $d=1$.
\end{enumerate}

\subsection{Comparison with the Riemann zeros}
\label{subsec:zeros_control}

To verify that the NEE phase is not an artifact of the zeta spectrum,
we evaluated the spectral form factor directly on the Riemann zeros.

Using the first $10\,000$ zeros and the standard ramp window $t \in [0.5, 4.0]$, block-averaged SFF yields
\begin{equation}
\gamma_{\text{zeros}} = 1.1159 \pm 0.0294 \quad (R^2 = 0.78),
\end{equation}
consistent with the ergodic GUE prediction $\gamma = 1.0$. Analyzing the full $100\,000$-zero dataset across $t \in [1, 10]$ gives $\gamma_{\text{zeros}} = 0.94 \pm 0.12$ ($R^2 = 0.21$), with a 10-chunk jackknife confirming stability ($\gamma = 0.95 \pm 0.31$). Both tests confirm that the actual Riemann zeros are ergodic, establishing that the sub-diffusive ramp ($\gamma \approx 0.61$) is an intrinsic consequence of the sparse hopping constraint.

\subsection{Coupling-strength dependence of the collapse}
\label{sec:coupling_sweep}

To test whether the spatial collapse without $d=1$ can be overcome by increasing the coupling strength, we performed a parameter sweep over $\varepsilon \in [0.2, 4.0]\,\varepsilon_c$ at $N = 2048$ with and without first-neighbor hopping (\cref{fig:coupling_sweep}).

\begin{figure}[htbp]
\centering
\includegraphics[width=0.85\textwidth]{Fig_Epsilon_Sweep.pdf}
\caption{Fractal dimension $D_2$ as a function of relative coupling strength $\varepsilon/\varepsilon_c$ at $N=2048$. The collapse upon removing $d=1$ is maximized at the critical resonance $\varepsilon = \varepsilon_c$ (a sixfold suppression). Even at $\varepsilon = 4\varepsilon_c$, $D_2$ without $d=1$ remains suppressed, demonstrating topological pinning.}
\label{fig:coupling_sweep}
\end{figure}

At the critical coupling $\varepsilon = \varepsilon_c$, removing $d=1$ reduces the fractal dimension from $D_2 = 0.2835$ to $D_2 = 0.0480$ (a collapse ratio of $0.17$). This demonstrates topological protection: at resonance, long-range hopping cannot substitute for the missing nearest-neighbor channel because the 1D backbone fragments into decoupled clusters.

As $\varepsilon$ increases beyond $\varepsilon_c$, higher-order virtual transitions across long-range channels progressively raise $D_2$. The ratio $r(\varepsilon) \equiv D_2^{\text{no}\,d=1} / D_2^{\text{with}\,d=1}$ increases monotonically, reaching $r(4\varepsilon_c) = 0.87$. We note that the regime $\varepsilon \gtrsim 2\varepsilon_c$ lies outside the calibration window of the PRBM model ($\varepsilon = \varepsilon_c$), so this saturation should be interpreted as a qualitative trend rather than a quantitatively reliable feature of the model; the topological protection at $\varepsilon = \varepsilon_c$ is the physically relevant regime. Even at $\varepsilon = 4\varepsilon_c$, $D_2$ without $d=1$ ($0.384$) remains substantially below the unconstrained value ($0.442$), so the collapse is never fully healed.

\section{Mechanized Formal Verification in Lean 4}
\label{sec:lean4_verification}

The algebraic, topological, and connectivity constraints governing the hopping mask are certified in the \textsf{Lean 4} Interactive Theorem Prover (v4.34.1) backed by \textsf{Mathlib4}. All formal proofs are 100\% \texttt{sorry}-free, verified by the \textsf{Lean 4} kernel at Level N1, and archived on Zenodo. The certified suite comprises 13 decoupled modules (23 theorems and 11 definitions, 34 total verified statements). We distinguish below between \emph{arithmetic} certificates (elementary number-theoretic identities specific to the mask structure) and \emph{physical} certificates (statements that connect the mask to quantitative properties of the interaction graph or the spectral dimension).

The certified statements fall into two categories: elementary arithmetic
identities specific to the coprime-to-6 mask structure, and quantitative
bounds that connect the mask to the interaction graph or to the fractal
dimension. Their mechanical verification guarantees that the algebraic
backbone of the paper contains no logical gaps that could be obscured by
notation or convention.

\subsection*{Arithmetic certificates}

\begin{enumerate}
    \item \textbf{Topological wheel channel structure} (\texttt{WheelChannelStructure.lean}):
    \begin{equation}
    \forall d \in \N, \quad \gcd(d, 6) = 1 \implies (d \bmod 6 = 1) \lor (d \bmod 6 = 5).
    \end{equation}

    \item \textbf{Unit cell measure} (\texttt{UnitCellMeasure.lean}):
    \begin{equation}
    \sum_{d=1}^{6} \Xi(d) = 2 \quad \implies \quad \varphi(6)/6 = 1/3.
    \end{equation}

    \item \textbf{Chiral phase absorption} (\texttt{ChiralPhaseAbsorption.lean}):
    \begin{equation}
    \forall d \in \N, \quad \gcd(d, 6) = 1 \implies \exists k \in \N, \quad (d \cdot k) \bmod 6 = 1.
    \end{equation}

    \item \textbf{Kato--Rellich window} (\texttt{KatoRellichRegime.lean}):
    \begin{equation}
    \nu = 3/4 \implies (2\nu > 1) \land (\nu < 1).
    \end{equation}

    \item \textbf{Local resonance coupling} (\texttt{LocalResonanceCoupling.lean}):
    \begin{equation}
    D = 2\pi \land 2\epsilon^2 = D^2 \implies \epsilon = \pi\sqrt{2}.
    \end{equation}

    \item \textbf{Channel density contraction} (\texttt{ChannelDensityContraction.lean}):
    \begin{equation}
    4/15 < 1/3 < 1/2 < 1.
    \end{equation}

    \item \textbf{Non-ergodic phase bound} (\texttt{NonErgodicPhaseBound.lean}): $1/4 < 1$.

    \item \textbf{Polyphase stopband annihilation} (\texttt{PolyphaseFilterZ6.lean}):
    \begin{equation}
    \gcd(d, 6) \neq 1 \implies d(d-2)(d-3)(d-4) \equiv 0 \pmod 6.
    \end{equation}
\end{enumerate}

\subsection*{Physical certificates}

\begin{enumerate}
    \setcounter{enumi}{8}
    \item \textbf{First-neighbor connectivity guarantee} (\texttt{FirstNeighborPreservation.lean}):
    \begin{equation}
    \Xi(1) = 1 \implies \Xi(1) \cdot w_1 = 2, \quad \text{and} \quad \Xi(1) = 0 \implies \Xi(1) \cdot w_1 = 0,
    \end{equation}
    certifying the guaranteed local interaction weight contribution $w_1 = 2$ under first-neighbor preservation (\cref{lem:degree_lower_bound}).

    \item \textbf{Dihedral quantization} (\texttt{DihedralQuantization.lean}): certifies that the chiral channel $d=5$ acts strictly as the reflection involution ($5x = -x \pmod 6$), bounding the quantum percolation threshold $k_{\min}$.

    \item \textbf{Bipartite quantum interference limit} (\texttt{BipartiteInterferenceBound.lean}): under chiral equiprobability $c_1 = c_5 = 1/2$, we have $c_1 \cdot c_5 = 1/4$, giving the upper bound $D_2 \le 1/4$.

    \item \textbf{Combinatorial degree bound \& totient density} (\texttt{ExpectedDegreeBound.lean}): certifies the first-neighbor local degree bound ($\langle \deg \rangle_{\text{local}} \ge 2$) and evaluates the coprime-to-6 channel density $\varphi(6)/6 = 1/3$.

    \item \textbf{Fractal dimension upper bound} (\texttt{FractalDimensionUpperBound.lean}): formalizes the discrete Jensen's inequality for the IPR ($p_1^2 + p_5^2 \ge \rho^2/2$), certifying the upper bound $D_2 \le \varphi(m)/m$ and its strict evaluation $D_2 \le 1/3$ for $m=6$; combined with module~11, it yields the sharpened bound $D_2 \le 1/4$ under chiral equiprobability.
\end{enumerate}

The execution of the verification suite is fully automated via Lake in a reproducible cloud kernel environment.

\section{Discussion}
\label{sec:discussion}

The central finding of this work is that the NEE phase in modular PRBM is not a consequence of the arithmetic structure of the mask, but of the preservation of short-range hopping. The $\Zmod$ coprime mask, the Möbius mask, and random masks with matched short-range support are spectrally equivalent. Removing $d = 1$ from the coprime mask collapses the NEE phase entirely.

This has several implications:

\textbf{(i) The arithmetic origin of the mask is a motivation, not a mechanism.} The noncommutative geometry approach to the Standard Model motivates the coprime-to-6 mask, but does not imply a distinctive spectral signature. The spectral properties follow from the mask's channel density and the preservation of short-range hopping, both of which are shared by simpler alternative masks.

\textbf{(ii) The M\"obius function exhibits no autonomous spectral signature beyond short-range variance dominance.}
The apparent match between the M\"obius mask and a random Bernoulli mask
at matched density ($\rho = 6/\pi^2$) is physically driven by power-law
variance concentration rather than asymptotic distribution identity. Any
coherent phase correlations from deep arithmetic channels ($d > 5$) have
amplitudes well below the finite-size GUE statistical noise floor, and
are therefore unresolvable at the accessible system sizes. Consequently,
the number-theoretic structure of $\mu(d)$—including square-free
multiplicativity and Sarnak-type orthogonality—acts effectively as
topological white noise, leaving no macroscopic imprint on the fractal
dimension $D_2$ or on the spectral form factor. This provides a mesoscopic
realization of the Möbius Randomness Principle within the PRBM universality
class.

\textbf{(iii) The hierarchy of bounds $D_2 \le 1/4 < 1/3 \le \varphi(m)/m$.} Three bounds coexist in the present framework, each with different hypotheses:
\begin{itemize}
    \item $D_2 \le \varphi(m)/m$ (generic, all masks, all moduli $m$), certified in \texttt{FractalDimensionUpperBound.lean}.
    \item $D_2 \le 1/3$ (specialization to $m=6$, using $\varphi(6)/6 = 1/3$), certified in \texttt{ExpectedDegreeBound.lean} and \texttt{FractalDimensionUpperBound.lean}.
    \item $D_2 \le 1/4$ (sharpened under chiral equiprobability $c_1 = c_5 = 1/2$), certified in\\ \texttt{BipartiteInterferenceBound.lean}.
\end{itemize}
The coprime-to-6 mask achieves $D_2 \approx 0.243$, close to the sharpest bound $1/4$. Random masks with the same density collapse to $D_2 \approx 0.056$, well below the bound, showing that the bound is not universally saturated. The physical interpretation is that the chiral equiprobability condition selects a maximum-entropy configuration of the active channels, which is compatible with the coprime-to-6 mask but not with generic random masks.

\textbf{(iv) The Anderson transition at $\nu > 1$ is modified by the mask in a regime-dependent way.} The mask reduces $\langle r \rangle$ for $\nu \in (1, 2.5)$ and becomes negligible for $\nu > 2.5$. This is consistent with the short-range nature of the mask: the coprime-6 constraint eliminates only $d \in \{0, 2, 3, 4 \pmod 6\}$, which are negligible at large $\nu$ where the hopping kernel is dominated by $d = 1, 5, 7$.

\textbf{(v) Topological protection at resonance vs.\ energetic percolation at strong coupling.} The sweep over $\varepsilon \in [0.2, 4.0]\,\varepsilon_c$ confirms that at the physical operating point $\varepsilon = \varepsilon_c$, the collapse without $d=1$ is strictly topological. Overcoming this localization requires a massive increase in coupling strength ($\varepsilon \ge 4\varepsilon_c$), demonstrating that nearest-neighbor bonds constitute the indispensable minimal spanning backbone for quantum percolation.

\textbf{(vi) Relation to the Hilbert--Pólya program.} Our operator is
not a candidate for the Hilbert--Pólya operator. It reproduces the
macroscopic Weyl density of the Riemann zeros ($R^2 > 0.999999$) and
lies within the same broad universality class of local GUE statistics,
but it does not reproduce the individual eigenvalues: after global
alignment, the median relative deviation between the model spectrum
and the true zeros is of order $10^{-4}$, with only a small fraction
of eigenvalues matching the zeros at the $10^{-5}$ level. Combined with
the opposite sign of the finite-size drift of the level spacing ratio,
this establishes that our operator provides a statistical realization
of the zeta spectrum in the PRBM universality class, not a mechanism
for the Riemann hypothesis. A detailed three-level diagnostic (density,
local statistics, eigenvalue identity) is reported in the companion
notebook.

\section{Conclusions}
\label{sec:conclusions}

We have constructed a hybrid deterministic-stochastic Hamiltonian combining an exact inverse Weyl law potential with a sparse hopping mask and PRBM long-range hopping. We have shown that the NEE phase in modular PRBM is entirely governed by the preservation of short-range hopping, not by the arithmetic structure of the mask. Random masks with matched density collapse when $d = 1$ is absent and restore when $d = 1$ is forced. The coprime-to-6 mask is spectrally equivalent to a random mask with $d = 1$ preserved, with $D_2(\Xi_6) = 0.2882 \pm 0.0014$ and $D_2(\Xi_6 \setminus \{d=1\}) = 0.0448 \pm 0.0009$, a collapse ratio of $\sim 0.16$. The M\"obius mask is spectrally equivalent to a random mask sharing its first four short-range channels $d \in \{1, 2, 3, 5\}$, governed by power-law coupling variance dominance ($86.05\%$ local weight) and the swamping of deep arithmetic channels by the GUE statistical noise floor ($\sigma_{\text{GUE}} \sim 1/\sqrt{N}$). Sweeping the coupling strength confirms that this collapse is topologically protected at the critical resonance $\varepsilon_c$.

First, the finite-size scaling collapse of the full IPR distribution
across a sixfold range in $N$ (\cref{subsec:fss_collapse}) provides the
decisive, model-independent proof of non-ergodicity: all data collapse
onto a scale-invariant Gaussian master profile $\mathcal{F}(x)$,
establishing $D_2^\infty = 0.113 \pm 0.003$ and rejecting
$D_2^\infty = 0$ at $39\sigma$ significance. Second, the Chalker
dynamical scaling relation $\beta = 2D_2$ is confirmed to $0.2\sigma$
between real-space wavepacket spreading and the bulk fractal dimension
(\cref{sec:gpu_diag}). Third, information-theoretic model selection
(\cref{subsec:aic_bic_selection}) demonstrates that under the standard
PRBM scaling ansatz ($D_2^\infty + a/\ln N$), the asymptotic fractal
dimension remains strictly positive at $81\sigma$ from zero
($D_2^\infty = 0.1225 \pm 0.0015$), while more flexible multiparameter
models suffer from severe variance inflation without exhibiting any
systematic drift of the underlying distribution toward localization.

The bound $D_2 \le \varphi(m)/m$ is certified rigorously in Lean~4 for
all modular masks, sharpened to $D_2 \le 1/4$ under chiral equiprobability,
and a conditional lower bound is derived for masks preserving first-neighbor
hopping with globally connected interaction graphs.

The measured Riemann zeros are ergodic ($\gamma_{\text{zeros}} \approx
0.94 \pm 0.12$ over $100\,000$ zeros), confirming that the NEE phase is
intrinsic to the sparse constraint. By uniting large-scale exact numerical
diagonalizations with kernel-certified Lean~4 formal proofs, this work
establishes an analytically tractable and mechanically verified model of
sparse-hopping quantum chaos.

\section*{Acknowledgments}
The author conducted this research independently and acknowledges the
open-source communities behind Lean~4 and Mathlib4 (formal verification),
PyTorch and SciPy (numerical diagonalization and statistical analysis),
Google Colab (GPU-accelerated experimental environment), and Google Gemini
and DeepSeek (AI-assisted language editing and critical review, see the
disclosure below). This research received no external funding.

\section*{Author Contributions (CRediT Taxonomy)}
\textbf{José Ignacio Peinador Sala:} Conceptualization, Methodology,
Software, Formal Analysis, Formal Verification, Investigation, Data
Curation, Validation, Writing -- Original Draft, Writing -- Review \&
Editing, Visualization, Project Administration.

\section*{Ethics Statement}
Ethical approval was not required for this study, as the research is
theoretical and numerical in nature and involves no human participants,
human data, tissue samples, or animal subjects.

\section*{Data and Code Availability}
\begin{itemize}
    \item \textbf{GitHub Repository:} \url{https://github.com/NachoPeinador/The-Riemann-GUE-Hamiltonian}
    \item \textbf{Archival Artifact (Zenodo DOI):} \url{https://doi.org/10.5281/zenodo.19284510}
\end{itemize}
The formal verification suite comprises 13 Level N1 modules certified by
the Lean~4 kernel (version 4.34.1 with Mathlib4), with 23 theorems and 11
definitions and zero \texttt{sorry} assertions.
SHA-256: 01ab66c630ecfa1be9fa9d82502518f8ccc14478b276ba040835f764ebcae6c5.
The numerical experiments
are reproducible from the companion Google Colab notebook archived in the
same Zenodo record. The notebook generates all figures and tables of this
manuscript from a fixed random seed (\texttt{SEED = 42}) and uses
double-precision (\texttt{complex128}) arithmetic throughout, as described
in the main text.

\section*{Licensing and Open Access}
This manuscript and its proofs are distributed under
\href{https://creativecommons.org/licenses/by/4.0/}{CC BY 4.0}. Source
code, including the Lean~4 verification suite and the numerical notebook,
is released under the MIT license.

\section*{Competing Interests}
The author declares no competing interests, financial or otherwise.

\section*{Declaration of Generative AI and AI-Assisted Technologies}
In the interest of transparency, the author discloses the use of generative
AI tools in the preparation of this manuscript. Google Gemini and DeepSeek
were employed exclusively for three auxiliary purposes: (i) English language
editing and stylistic refinement of the manuscript text; (ii) structural
review and critical feedback on the logical flow of the argument, acting in
an adversarial ``red-team'' capacity to stress-test the theoretical claims;
and (iii) assistance in the implementation and debugging of the
reproducibility notebook and LaTeX source files. The generative models were
not used to produce any original scientific content, including mathematical
derivations, numerical results, theoretical postulates, or conjectures.
All substantive intellectual contributions---the formulation of the model,
the design of the numerical experiments, the derivation of the bounds, the
formal verification suite, and the interpretation of the results---are the
author's own work. The author has reviewed, verified, and validated every
AI-assisted textual and code output, and takes full responsibility for the
final content of the manuscript.

\section*{Suggested Citation}
J.~I.~Peinador~Sala, \textit{First-neighbor connectivity governs the
non-ergodic extended phase in power-law random banded matrices with sparse
hopping constraints}, Zenodo Archive (2026).
\url{https://doi.org/10.5281/zenodo.19284510}

\appendix

\section{Derivation of \texorpdfstring{$\epsilon_c$}{epsilon\_c} and the fractional ramp}
\label{app:ramp}

\subsection{Resonance condition for \texorpdfstring{$\epsilon_c$}{epsilon\_c}}
The local level spacing is $D = 2\pi$. For hybridization, $2\sigma_V^2 = D^2$ with $\sigma_V^2 = \epsilon^2$, giving $\epsilon_c = \pi\sqrt{2}$. This derivation from the resonance constraint is certified in the Lean~4 module \texttt{LocalResonanceCoupling.lean}.

\subsection{Altshuler--Shklovskii scaling and the SFF ramp}
For a PRBM in the NEE phase, $\Sigma^2(\Delta E) \sim (\Delta E)^{D_2}$, so the ramp behaves as $K_{\text{conn}}(t) \propto t^{\gamma}$ with $\gamma = D_2 + \eta$. With $D_2 \approx 0.243$ and $\eta \approx 0.365$, we obtain $\gamma \approx 0.608$, consistent with the numerical fit.

\section{Formal upper bound \texorpdfstring{$D_2 \le \varphi(m)/m$}{D2 <= phi(m)/m}}
\label{app:bound}

The theoretical upper bound for the fractal dimension $D_2$ is fundamentally constrained by the active channel density $\rho = \varphi(m)/m$. This relationship is formalized in the Lean~4 module \texttt{FractalDimensionUpperBound.lean} via the discrete Jensen's inequality for the Inverse Participation Ratio (IPR).

For a state strictly supported on the active chiral channels ($d \equiv 1$ and $d \equiv 5 \pmod 6$), let $p_1$ and $p_5$ represent the fractional spectral weights of these sublattices. Under the chiral equiprobability assumption $p_1 = p_5$, with total weight $p_1 + p_5 = \rho$, the harmonic mean core gives
\begin{equation}
\IPR \propto p_1^2 + p_5^2 \ge \frac{(p_1 + p_5)^2}{2} = \frac{\rho^2}{2}.
\end{equation}
Because maximizing spatial extent (and thus maximizing $D_2$) requires minimizing the IPR, the fractal dimension is bounded by the density of available quantum channels.

\paragraph{From channel IPR to global $D_2$.}
The inequality above bounds the IPR restricted to the active chiral
sublattice. To lift it to the full state, we assume that the wavefunction
is supported predominantly on the active channels: inactive channels have
vanishing hopping amplitude and are therefore dynamically decoupled.
Under this assumption the global IPR coincides with the sublattice IPR,
and the bound applies to $D_2$ directly. Relaxing the assumption (i.e.,
allowing leakage onto inactive channels) can only increase the IPR and
therefore can only decrease $D_2$; the bound is thus robust.

Combining the chiral equiprobability assumption with the totient density $\rho = \varphi(6)/6 = 1/3$ (certified in \texttt{ExpectedDegreeBound.lean}), we obtain the hierarchy
\begin{equation}
D_2 \le 1/4 \quad (\text{chiral equiprobability}),
\end{equation}
which is strictly sharper than the generic bound
\begin{equation}
D_2 \le 1/3 \quad (\text{specialization to } m=6),
\end{equation}
which in turn is a specialization of
\begin{equation}
D_2 \le \varphi(m)/m \quad (\text{all moduli } m).
\end{equation}

The sharpest certified bound $D_2 \le 1/4$ is consistent with our empirical GPU bulk median measurements $\langle D_2 \rangle_{\text{median}} \approx 0.25$ (\cref{sec:gpu_diag}). We emphasize that this is a consistency check rather than a demonstration of equality.

\paragraph{Rigorous derivation: from modular cosets to global $D_2$.}
To formalize the connection between the discrete chiral bound and the spatial $\IPR = \sum_j |\psi(j)|^4$, we partition the lattice $\Lambda = \{1, \dots, N\}$ into the six residue classes (cosets) of $\Zmod$:
\begin{equation}
\mathcal{S}_r := \{j \in \Lambda \mid j \equiv r \pmod 6\}, \quad r \in \{0, 1, \dots, 5\},
\end{equation}
each containing $N_r = N/6 + O(1)$ sites. Let $W_r := \sum_{j \in \mathcal{S}_r} |\psi(j)|^2$ denote the total probability weight residing on coset $r$, satisfying $\sum_{r=0}^5 W_r = 1$. Applying Cauchy--Schwarz within each coset yields
\begin{equation}
\sum_{j \in \mathcal{S}_r} |\psi(j)|^4 \ge \frac{1}{N_r}\left( \sum_{j \in \mathcal{S}_r} |\psi(j)|^2 \right)^2 = \frac{6}{N} W_r^2.
\end{equation}
Summing over all six cosets establishes the exact relation
\begin{equation}
\IPR \ge \frac{6}{N} \sum_{r=0}^5 W_r^2.
\label{eq:ipr_coset_bound}
\end{equation}
Because the hopping operator $\hat{V}$ only permits steps with
$d \equiv \pm 1 \pmod 6$, quantum transitions strictly couple adjacent
cosets, $r \leftrightarrow (r \pm 1) \pmod 6$, forming a discrete
6-cycle. Direct transitions on stopband channels
($d \equiv 0, 2, 3, 4 \pmod 6$) vanish identically at first order
($V_{ij} = 0$). It remains to verify that second-order virtual processes
do not reintroduce an effective coupling between non-adjacent cosets
at the accessible system sizes.

Second-order virtual processes mediating transitions across inactive intervals (e.g., $i \to i+2$ via intermediate state $k = i+1$) generate an effective coupling scaling via resolvent perturbation theory as
\begin{equation}
V_{\text{eff}}^{(2)}(i, i+2) = \sum_{k} \frac{V_{ik} V_{k, i+2}}{E - E_k} \propto \frac{\epsilon^2}{\delta E} \, d_1^{-\nu} d_2^{-\nu}.
\end{equation}
Since $\nu = 0.75 > 0.5$, the total off-shell leakage probability onto non-chiral multi-hop pathways is bounded by $\mathcal{P}_{\text{leak}}(N) \le C N^{1 - 2\nu} = C N^{-0.5} \to 0$ as $N \to \infty$. Consequently, the chiral restriction remains non-perturbative in the thermodynamic limit, certifying that the Jensen upper bound $D_2 \le 1/4$ holds asymptotically without preasymptotic leakage corrections.

\paragraph{Empirical validation of the chiral equiprobability axiom.}
In the Lean~4 proof of the bipartite upper bound $D_2 \le 1/4$ (\texttt{BipartiteInterferenceBound.lean}), the assumption of chiral equiprobability $c_1 = c_5 = 1/2$ is injected axiomatically. To ensure that this algebraic condition is not violated by spontaneous local mass fluctuations, we empirically computed the sublattice probability weights $W_r \equiv \sum_{j \equiv r \pmod 6} |\psi(j)|^2$ across all six residue classes of $\Zmod$ on GPU bulk eigenmodes up to $N = 12\,000$.

\begin{figure}[htbp]
\centering
\includegraphics[width=0.92\textwidth]{Fig_Chiral_Equiprobability_Audit.pdf}
\caption{\textbf{Empirical audit of the chiral equiprobability axiom.}
\textbf{(a)} Sublattice probability weights $\langle W_r \rangle$ across
all six residue classes of $\Zmod$ at $N = 12\,000$. The ensemble mean
satisfies exact equipartition ($\langle W_r \rangle \approx 1/6$) with
net chiral asymmetry $\langle W_1 - W_5 \rangle = -0.000048 \approx 0$.
\textbf{(b)} Mean absolute state-by-state chiral asymmetry
$\langle |W_1 - W_5| \rangle$ as a function of $N$, showing that local
mass fluctuations average out in the thermodynamic limit.}
\label{fig:chiral_audit}
\end{figure}

As shown in \cref{fig:chiral_audit}(a), the ensemble-averaged weights across all six cosets satisfy exact equipartition to four decimal places:
\begin{equation}
\langle W_r \rangle = 0.1667 \approx 1/6 \quad \forall r \in \{0, 1, 2, 3, 4, 5\}.
\end{equation}
The net chiral imbalance vanishes identically: $\langle W_1 - W_5 \rangle = -0.000048 \approx 0$. Conditioning on the active chiral subspace ($W_1 + W_5 = \rho = 1/3$), the effective relative weights are $c_1 = c_5 = 0.5000$. Furthermore, \cref{fig:chiral_audit}(b) demonstrates that the mean absolute state-by-state asymmetry $\langle |W_1 - W_5| \rangle$ decays monotonically from $0.1246$ at $N = 2000$ to $0.1116$ at $N = 12\,000$. This confirms that local deviations represent finite-size Gaussian fluctuations, validating the thermodynamic legitimacy of the Lean~4 certified bound.

\section{Alternative masks: technical details}
\label{app:altmasks}

The Möbius function $\mu(d)$ is defined as $0$ for non-square-free integers and $(-1)^k$ for integers with $k$ distinct prime factors. Its density of nonzero values is $6/\pi^2 \approx 0.608$. The Liouville function $\lambda(d) = (-1)^{\Omega(d)}$ (where $\Omega$ counts prime factors with multiplicity) has density $1/2$ for $\lambda(d) = -1$.

Random masks are generated by independent Bernoulli sampling with probability $\rho$. Forced-channel masks set $\Xi(d) = 1$ for $d \in D_{\text{forced}}$, and Bernoulli($\rho$) otherwise, without re-normalization.

\section{Lower bound on \texorpdfstring{$D_2$}{D2} from local connectivity preservation}
\label{app:lower_bound}

In this appendix we establish a conditional lower bound on the fractal dimension $D_2$ that holds whenever the hopping mask preserves first-neighbor connectivity \emph{and} the resulting interaction graph remains globally connected. The bound is stated in two parts: a rigorous theorem relating the mask to the expected degree of the interaction graph, and a scaling relation connecting the degree to $D_2$ under the standard PRBM ansatz.

\subsection{Expected degree under local connectivity preservation}
\label{subsec:expected_degree}

Let $\Xi : \N \to \{0, 1\}$ be a mask with $\Xi(1) = 1$ (first-neighbor hopping preserved). Define the expected degree of the interaction graph at system size $N$ by
\begin{equation}
\langle \deg \rangle_N(\Xi, \nu) := \sum_{\substack{1 \le d \le N-1 \\ \Xi(d) = 1}} w_d, \quad w_d := 2 \, d^{-2\nu},
\label{eq:expected_degree_def}
\end{equation}
where the factor $2$ accounts for the two spatial neighbors (left and right) at distance $d$ from a site deep in the bulk of the chain. For boundary sites, this factor is asymmetric and locally reduced; such finite-size boundary effects vanish in the thermodynamic limit $N \to \infty$, rendering the bulk approximation asymptotically exact.

\begin{lemma}[Lower bound on the expected degree]
\label{lem:degree_lower_bound}
For any mask $\Xi$ with $\Xi(1) = 1$, and for any $\nu \in (0.5, 1)$,
\begin{equation}
\langle \deg \rangle_N(\Xi, \nu) \geq 2 + \frac{\varphi(m)}{m} \cdot \frac{N^{1-2\nu}}{1 - 2\nu} + O(1),
\label{eq:degree_lower}
\end{equation}
where the first term is the guaranteed first-neighbor contribution and the second term comes from long-range channels.
\end{lemma}

\begin{proof}
The contribution of $d = 1$ is exactly $w_1 = 2$, independently of $N$ and $\nu$. This base contribution is certified in the Lean~4 module \texttt{FirstNeighborPreservation.lean}. For $d \geq 2$, bounding the sum below by the coprime channel fraction $\varphi(m)/m$ times the full sum from $d = 2$ to $N-1$ (asymptotically equivalent to the full sum), we obtain \cref{eq:degree_lower}.
\end{proof}

\begin{corollary}[Quantitative bound at $\nu = 0.75$]
\label{cor:degree_nu075}
At $\nu = 0.75$, $1 - 2\nu = -0.5$, so
\begin{equation}
\langle \deg \rangle_N \geq 2 + 2 \cdot \frac{\varphi(m)}{m} \cdot N^{0.5} + O(1).
\label{eq:degree_at_075}
\end{equation}
For the coprime-to-6 mask ($m = 6$, $\varphi(6)/6 = 1/3$) and $N = 1500$:
\begin{equation}
\langle \deg \rangle_{1500} \geq 2 + \frac{2}{3} \cdot \sqrt{1500} \approx 2 + 25.8 = 27.8.
\end{equation}
\end{corollary}

\subsection{From expected degree to fractal dimension: an effective \texorpdfstring{$k_{\min}$}{k\_min}}
\label{subsec:degree_to_fractal}

The bridge connecting the expected degree $\langle \deg \rangle_N$ of the interaction graph to the fractal dimension $D_2$ rests on the PRBM finite-size scaling hypothesis \cite{Mirlin1996}, which gives
\begin{equation}
D_2 \geq \frac{\log(k_N / k_{\min})}{\log N} + o(1),
\label{eq:degree_fractal_revised}
\end{equation}
where $k_{\min} > 0$ is the minimum connectivity threshold required to sustain quantum percolation.

In our $\Zmod$ masked model, this threshold is structurally constrained by the algebraic properties of the modular ring. The active channels are restricted to the identity $d=1$ and its chiral reflection $d=5$ (since $5 \equiv -1 \pmod 6$), as certified in \texttt{DihedralQuantization.lean}; intermediate channels are suppressed by the topological stopband (\texttt{PolyphaseFilterZ6.lean}). Because a nearest-neighbor transition $i \to i+1$ in the absence of the $d=1$ bond would require a second-order virtual process through the $d=5$ channel, subject to severe phase detuning from the diagonal potential $E_k$, the base coordination number must guarantee the $d=1$ forward channel to prevent the 1D backbone from fragmenting into isolated clusters. This provides a topological estimate for the localization threshold,
\begin{equation}
k_{\min} \in [1.0, 1.7],
\end{equation}
consistent with the empirical value $k_{\min} \approx 1.67$ extracted from exact GPU diagonalizations (\cref{subsec:numerical_consistency}).

\subsection{Combined lower bound}
\label{subsec:combined_bound}

Combining \cref{lem:degree_lower_bound} with the scaling ansatz \cref{eq:degree_fractal_revised}, we obtain the following conditional lower bound on $D_2$.

\begin{theorem}[Conditional lower bound on $D_2$]
\label{thm:lower_bound_D2}
\emph{Assume} (i) the standard PRBM scaling relation $D_2 \geq \log(k_N/k_{\min})/\log N$, and (ii) that the interaction graph induced by the mask is globally connected. \emph{Then} for any mask $\Xi$ with $\Xi(1) = 1$ and $\nu = 0.75$,
\begin{equation}
D_2 \geq \frac{\log\!\left(2 + 2 \cdot \frac{\varphi(m)}{m} \cdot N^{0.5} + O(1)\right) - \log k_{\min}}{\log N}.
\label{eq:lower_bound_D2}
\end{equation}
For large $N$ and $m = 6$,
\begin{equation}
D_2 \geq \frac{1}{2} + \frac{\log(\varphi(6)/6) - \log k_{\min}}{\log N} + o(1/\log N).
\label{eq:lower_bound_asymptotic}
\end{equation}
\end{theorem}

\begin{proof}
Substitute \cref{eq:degree_at_075} into \cref{eq:degree_fractal_revised} and simplify. For $N \to \infty$, the first term inside the log is dominated by the $N^{0.5}$ contribution, giving \cref{eq:lower_bound_asymptotic}.
\end{proof}

The asymptotic form \cref{eq:lower_bound_asymptotic} shows that the lower bound is of the form
\begin{equation}
D_2 \geq \frac{1}{2} + \frac{c(m)}{\log N}, \quad c(m) := \log(\varphi(m)/m) - \log k_{\min},
\end{equation}
where the constant $c(m)$ depends on the mask only through the coprime density fraction $\varphi(m)/m$.

\subsection{Numerical consistency check}
\label{subsec:numerical_consistency}

To fix the order of magnitude of $k_{\min}$, we use the numerical values at $N = 1500$, $\nu = 0.75$ for the coprime-to-6 mask. The observed $D_2$ is $0.2882 \pm 0.0014$. Solving for $k_{\min}$ from \cref{eq:lower_bound_D2}:
\begin{equation}
k_{\min} \leq \frac{1 + \frac{1}{3} \sqrt{1500}}{N^{D_2^{\text{obs}}}} \approx \frac{13.9}{1500^{0.2882}} \approx \frac{13.9}{8.3} \approx 1.67.
\end{equation}
Since $k_{\min} \geq 1$ by definition, the numerical data is consistent with $k_{\min} \in [1, 1.7]$.

To verify the role of the local connectivity condition, we apply the same lower bound to a random mask with $\rho = 1/3$ and no forced $d = 1$. In this case, $\Xi(1) = 0$ with probability $2/3$, so the guaranteed first-neighbor contribution disappears. The lower bound of \cref{lem:degree_lower_bound} no longer applies; only the general bound holds, giving
\begin{equation}
\langle \deg \rangle_N^{\text{rand}} \sim \frac{2}{9} N^{0.5}.
\end{equation}
Inserting into \cref{eq:degree_fractal_revised} and taking $k_{\min} = 1$,
\begin{equation}
D_2^{\text{rand}} \geq 0.5 - \frac{\log(9/2)}{\log N} \approx 0.5 - 0.21 = 0.29 \quad \text{at } N = 1500.
\end{equation}
But the observed value is $D_2^{\text{rand}} = 0.0556 \pm 0.0007$. The lower bound is therefore \emph{violated} by the numerical data. This violation is the signature of the global-connectivity hypothesis in \cref{thm:lower_bound_D2} failing: random masks without $d=1$ fragment the graph, so the scaling ansatz \cref{eq:degree_fractal_revised}, which assumes a globally connected interaction graph, does not apply.

\subsection{Quantum percolation and cluster fragmentation upon removing \texorpdfstring{$d=1$}{d=1}}
\label{subsec:quantum_percolation}

The catastrophic collapse of the fractal dimension when first-neighbor hopping is removed is explained by the topological role of the $d=1$ bond in one spatial dimension.

\begin{proposition}[1D topological connectivity threshold]
\label{prop:connectivity_threshold}
Let $G = (V, E)$ be a one-dimensional lattice graph with hopping edges $(i, j)$ active whenever $|i-j| \in S_{\text{active}}$. If $1 \notin S_{\text{active}}$, then the shortest path distance $d_G(i, i+1)$ between adjacent sites in the spatial lattice scales as $\Theta(\min \{d \in S_{\text{active}}\})$.
\end{proposition}

\begin{proof}[Percolation analysis]
The spatial dimension $d=1$ plays a singular topological role: nearest-neighbor bonds ($d=1$) constitute the unique minimal spanning tree that guarantees global graph connectivity without spatial gaps.

\textbf{Step 1 (Creation of spatial bottlenecks).} If $\Xi(1) = 0$, any quantum transition between site $i$ and $i+1$ must proceed via higher-order indirect paths, e.g., $i \to i+5 \to i+1$. The effective hopping matrix element for this second-order virtual process scales, by perturbation theory, as
\begin{equation}
V_{\text{eff}}^{(2)}(i, i+1) \sim \frac{V_{i, i+5} V_{i+5, i+1}}{E_i - E_{i+5}} \propto d_1^{-\nu} d_2^{-\nu} \ll d^{-1},
\end{equation}
suffering from severe energy detuning due to the diagonal potential $E_k$.

\textbf{Step 2 (Quantum percolation failure).} Without $d=1$, the physical network fragments into a set of weakly coupled sparse clusters. Although the total degree $\langle \deg \rangle_N \sim \rho N^{1-2\nu}$ remains formally large due to long-range tails, the \emph{local spatial connectivity} vanishes. The system sits below the 1D quantum percolation threshold, causing the wavefunctions to localize within isolated spatial islands of size $\ell_{\text{cluster}} \sim 1/\rho$.

Consequently, the assumption of global wavefunction expansion underlying \cref{eq:degree_fractal_revised} fails, and the effective parameter $k_{\min} \to \infty$, correctly explaining the collapse to $D_2 \approx 0.05$.
\end{proof}

This percolation picture provides the physical interpretation of the conditional hypothesis in \cref{thm:lower_bound_D2}: the bound is a diagnostic of global connectivity, not a universal constraint on $D_2$. Masks that preserve $d=1$ but otherwise have low density may still violate the bound if they fail to maintain global connectivity; conversely, masks without $d=1$ but with sufficient long-range structure might in principle restore connectivity and satisfy the bound. For example, a mask that includes all distances $d \in \{1\} \cup \{2k : k \geq 1\}$
preserves first-neighbor connectivity and maintains dense long-range
support, and is expected to satisfy the bound; conversely, a mask that
excludes $d=1$ but includes all even distances is fragmented into parity
sectors and violates the global-connectivity hypothesis regardless of its
density. In the present study, the coprime-to-6 mask sits firmly in the first category (globally connected, bound satisfied), while generic random masks at $\rho = 1/3$ sit in the second (fragmented, bound violated).

\bibliographystyle{apsrev4-2}
\bibliography{references}

\end{document}